\documentclass[11pt]{amsart}
\usepackage[margin=30mm]{geometry}
\usepackage{amsmath,amssymb,amsthm,mathtools,microtype,hyperref}
\hypersetup{colorlinks=true,linkcolor=blue,citecolor=blue,urlcolor=blue}
\title[A Proof Candidate for Carlet's Conjecture]{A Complete Proof Candidate for Carlet's Conjecture\\[0.4em]\large Reciprocal Compactifications, Quadratic Complete Intersections, and Explicit Betti Bounds}
\author{Wang Quanzhi}
\author{ChatGPT}
\address{AI-assisted editorial contribution}
\email{mrhwrn120504@gmail.com}
\date{November 2026; revised from the August 2026 manuscript}
\begin{document}
\begin{abstract}
This paper organizes a complete proof candidate for the zero-sum subspace form of Carlet's conjecture. For odd $n$, the target is that for every $3\le k\le n-3$ there exists a $k$-dimensional $\mathbb F_2$-subspace $E\subset\mathbb F_{2^n}$ such that $\sum_{x\in E\setminus\{0\}}x^{-1}=0$. The proposed route converts the zero-sum condition into the vanishing of the $X^2$ coefficient of a subspace polynomial, realizes it as a Cartier section on a Pink--Schieder reciprocal compactification, constructs an explicit quadratic model for the full-rank stratum, and derives point-count estimates from stratified cohomology and finite certificates. Three explicit certificates for $(n,k)=(29,13),(29,14),(31,14)$ have been independently checked by exact computation. The general reciprocal-geometric argument in Sections 4--10 has not received a complete independent line-by-line review; accordingly, the main statement is presented as a proof candidate rather than as a certified theorem.
\end{abstract}
\maketitle
\begin{center}\fbox{\begin{minipage}{0.92\linewidth}\small\textbf{Manuscript status.} The three finite-field certificates have been independently checked exactly. This does not certify the general geometric argument. The main proof candidate remains subject to independent specialist review and has not been formalized in Lean.\end{minipage}}\end{center}
\section*{Revision note (November 2026)}
This release removes the earlier word-processing cover material and presents the manuscript with a journal-style mathematical front matter. It adds clarifications on (i) defining the formal product $L_j(T)$ over an arbitrary $\mathbb F_2$-algebra before using any linear-independence conclusion; (ii) the vanishing range of compactly supported cohomology used in the total Betti bound; and (iii) the Frobenius-compatible identification of the local intersection complex with the constant sheaf, including the Tate twist in duality. The complete arguments are recorded in the Chinese edition's corresponding sections; this English edition's source proof text follows the supplied English candidate manuscript.

On 5 October 2026, Kaimin Cheng posted the preprint \emph{Counting zero-sum subspaces for the multiplicative inverse function} (arXiv:2610.06445), whose abstract claims a proof of the full conjecture by a counting and Fourier-analytic method. We have not audited that proof and make no priority determination. Its existence is relevant to the current literature and is recorded here without assessing its correctness.

\noindent\textbf{Keywords:} Carlet conjecture; finite fields; multiplicative inverse; subspace polynomials; reciprocal planes; intersection cohomology; Betti bounds.
\newpage
\section*{Supplementary mathematical clarifications}
\textbf{1. The reciprocal polynomial over an arbitrary base ring.} Let $A$ be any commutative $\mathbb F_2$-algebra and let $x_1,\ldots,x_j\in A$. Define the formal polynomial
\[
L_j(T)=\prod_{\epsilon\in\mathbb F_2^j}\left(T+\sum_{i=1}^j\epsilon_i x_i\right)\in A[T].
\]
This product and its recursion are defined before any assertion that the $x_i$ are linearly independent. In the reverse construction of Theorem 4.2, the unit conditions first imply that every nonzero binary linear combination is a unit; only then does the tuple define a point of the Drinfeld open.

\textbf{2. The full degree range in the compact-support Betti sum.} The variety $X_r$ is smooth, affine, and pure of dimension $d=r-2$. Artin vanishing gives $H^j(X_r)=0$ for $j>d$, and Poincar\'e duality gives $H_c^i(X_r)=0$ for $i<d$. The dimension bound gives $H_c^i(X_r)=0$ for $i>2d$. Thus the sum in Lemma 7.1 includes every potentially nonzero compactly supported cohomology group; the case $r=2$ has $d=0$ and is covered as well.

\textbf{3. Frobenius structure and duality.} The identification of the local intersection complex with the shifted constant sheaf must be taken compatibly with the structure over $\mathbb F_2$, hence with Frobenius. For a proper rationally smooth $m$-dimensional variety, the duality used in (6.2) is $H^i\cong H^{2m-i,\vee}(-m)$, with its Tate twist retained. These clarifications make explicit the Frobenius representation used by the trace formula.

\textbf{Related work.} Kaimin Cheng's preprint, submitted to arXiv on 5 October 2026, claims a proof of the full conjecture by a counting recurrence and Fourier analysis (arXiv:2610.06445v1). We have not audited that proof and make no priority determination.
\end{document}
